\documentclass[10pt]{article}
\usepackage[T1]{fontenc}
\usepackage[utf8]{inputenc}
\usepackage[french]{babel}
\usepackage{amsmath,amssymb,mathrsfs,mathpazo}
\usepackage{geometry}
\geometry{a4paper,margin=2cm}
\usepackage{graphicx,array,tabularx,tabularray,colortbl,multirow,multicol,booktabs,enumitem,pifont,ulem,fancybox,diagbox,cancel}
\usepackage{pstricks,pstricks-add,pst-plot,pst-tree,pst-node,pst-text,pst-func,pst-math,pst-bspline,pst-3dplot,pst-eucl,pst-solides3d}
\let\pscalculate\undefined
\usepackage{pst-bezier}
\usepackage{tikz,pgfplots,tkz-tab}
\usetikzlibrary{arrows,arrows.meta,patterns,calc,positioning,trees,decorations.markings,shapes.geometric,shapes.misc,shapes.arrows,intersections,decorations.pathmorphing,decorations.pathreplacing,backgrounds,shadows,babel}
\pgfplotsset{compat=1.18}
\usepgfplotslibrary{fillbetween}
\usepackage[np]{numprint}
\usepackage{listings,makecell,pas-tableur,twoopt}
\let\barre\undefined
\usepackage[pstricks]{tablvar}
\definecolor{darkgreen}{RGB}{0,100,0}
\providecommand{\red}{\color{red}}
\providecommand{\blue}{\color{blue}}
\providecommand{\green}{\color{green}}
\providecommand{\black}{\color{black}}
\providecommand{\white}{\color{white}}
\providecommand{\gray}{\color{gray}}
\providecommand{\cyan}{\color{cyan}}
\providecommand{\magenta}{\color{magenta}}
\providecommand{\yellow}{\color{yellow}}
\providecommand{\violet}{\color{violet}}
\providecommand{\orange}{\color{orange}}
\providecommand{\euro}{\texteuro}
\providecommand{\EUR}{\texteuro}
\providecommand{\ieme}{\textsuperscript{e}}
\providecommand{\degres}{\ensuremath{^\circ}}
\providecommand{\degree}{\ensuremath{^\circ}}
\providecommand{\vv}[1]{\overrightarrow{#1}}
\providecommand{\Coord}[1]{#1}
\providecommand{\decosix}{}
\providecommand{\starredbullet}{\ensuremath{\bullet}}
\pagestyle{empty}
\setlength{\parindent}{0pt}
\usepackage{adjustbox}\DeclareUnicodeCharacter{2236}{:}\DeclareUnicodeCharacter{2190}{\ensuremath{\leftarrow}}
\providecommand{\R}{\mathbb{R}}
\providecommand{\N}{\mathbb{N}}
\providecommand{\Z}{\mathbb{Z}}
\providecommand{\Q}{\mathbb{Q}}
\providecommand{\D}{\mathbb{D}}
\providecommand{\C}{\mathbb{C}}
\providecommand{\e}{{\mathrm{e}}}
\providecommand{\vect}[1]{\overrightarrow{#1}}
\providecommand{\vectt}[1]{\overrightarrow{\mathrm{#1}}}
\providecommand{\vvt}[1]{\overrightarrow{\mathrm{#1}}}
\providecommand{\np}[2][]{{#2~#1}}
\providecommand{\vv}[1]{{\overrightarrow{#1}}}
\renewcommand{\d}{{\mathrm{d}}}
\providecommand{\degres}{{^\circ}}
\providecommand{\degre}{{^\circ}}
\providecommand{\blue}{\color{blue}}
\providecommand{\red}{\color{red}}
\providecommand{\enskip}{\quad}
\providecommand{\llbracket}{\lbrack\!\lbrack}
\providecommand{\rrbracket}{\rbrack\!\rbrack}
\providecommand{\euro}{{\text{€}}}
\providecommand{\no}{{\text{n°}}}
\providecommand{\interval}[2]{{[#1;#2]}}
\providecommand{\intervFO}[2]{{[#1;#2[}}
\providecommand{\dotfill}{\ldots}

\providecommand{\ds}{\displaystyle}
\providecommand{\textding}[1]{#1}
\providecommand{\Oij}{\(\mathrm{O};\vec{\imath},\vec{\jmath}\)}
\providecommand{\Oijk}{\(\mathrm{O};\vec{\imath},\vec{\jmath},\vec{k}\)}
\providecommand{\Ouv}{\(\mathrm{O};\vec{u},\vec{v}\)}
\providecommand{\intervFF}[2]{\ensuremath{[#1;#2]}}
\renewcommand{\euro}{\ifmmode\text{\texteuro}\else\texteuro\fi}
\DeclareUnicodeCharacter{00B0}{\ifmmode^{\circ}\else\textdegree\fi}
\DeclareUnicodeCharacter{20AC}{\ifmmode\text{\texteuro}\else\texteuro\fi}
% Énoncés : Éducation nationale. Transcription et corrigés : APMEP, Denis Vergès et contributeurs. Découpage et classement : CoopMaths, Sébastien Lozano, Éric Elter et contributeurs. Adaptation de la présentation : Maths Explorer.
% CC BY-SA : https://coopmaths.fr/www/about/
% https://coopmaths.fr/alea/static/e3c/2020/tex/e3c_2020_00_sujet29_3_cor.tex
\begin{document}
\noindent\textit{APMEP / CoopMaths. CC BY-SA. Présentation adaptée par Maths Explorer.}\par\medskip
\begin{center}
\psset{yunit=3cm,labelFontSize=\scriptstyle}
\begin{pspicture}(-0.75,-0.2)(11,1.1)
\multido{\n=0.0+0.2}{56}{\psline[linewidth=0.25pt,linecolor=lightgray](\n,0)(\n,1)}
\multido{\n=0.0+0.2}{6}{\psline[linewidth=0.25pt,linecolor=lightgray](0,\n)(11,\n)}
\multido{\n=0+1}{11}{\psline[linewidth=0.45pt](\n,0)(\n,1)}
\multido{\n=0+1}{2}{\psline[linewidth=0.45pt](0,\n)(10.8,\n)}
\psaxes[linewidth=1.25pt]{->}(0,0)(-0.1,0)(11,1)
\psaxes[linewidth=1.25pt](0,0)(0,0)(11,1)


\def\Func{2.71828 x 0.5 neg  mul exp  x mul  }
\psplot[plotpoints=3000,linewidth=1.25pt,linecolor=red]{0}{11}{\Func}


\psline[linewidth=1.25pt,linecolor=blue]{->}(2,0.74)(0,0.74) 
\psline[linewidth=1.25pt,linecolor=blue]{->}(2,0)(2,0.74)    
\uput[l](0,0.74){\blue $f(2) \approx 0,74$}

\uput[dr](6,0.5){\red $\mathcal{C}$}

\end{pspicture}
\end{center}

\subsection*{1.}

On a :
\(
f(4) = 4 \times \e^{-0,5 \times 4} = 4 \times \e^{-2} = \dfrac{4}{\e^2}.
\)

\subsection*{2.}

\(f\) est un produit de fonctions dérivables sur \([0\,;\,+\infty[\), donc sur cet intervalle :
\[
f'(t) = \e^{-0{,}5t} - 0{,}5t \e^{-0{,}5t} = \e^{-0{,}5t}(1 - 0{,}5t).
\]

\subsection*{3.}

On sait que, quel que soit le réel \(t\), \(\e^{-0,5t} > 0\), donc \(f'(t)\) a le signe de \(1 - 0,5t\) :
\begin{align*}
&1 - 0,5t > 0 \\
\iff & 1 > 0,5t \\
\iff & 2 > t \\
\iff & t < 2 
\end{align*}
   

\subsection*{4.}

On en déduit que la fonction est croissante sur \([0\,;\,2]\) de \(f(0) = 0\) à \(f(2) = 2 \e^{-1}\), puis décroissante sur \([2\,;\,+\infty[\) de \(f(2)\) à zéro.
	\begin{center}
	\begin{tikzpicture}[double distance=2pt]
		\tkzTabInit{$x$/1,$f'(x)$/1,$f(x)$/2}{$0$,$2$, $+\infty$}
		\tkzTabLine{,+,z, -}
		\tkzTabVar{-/$0$,+/ $2 \e^{-1}$ /, -/}
	\end{tikzpicture}
\end{center} 

\subsection*{5.}

La question précédente montre que \(f(2) = 2 \e^{-1}\) est le maximum de la fonction sur \([0\,;\,+\infty[\).

\(f(2) = 2 \e^{-1} \approx 0,736\), soit environ 0,74, ce que confirme le graphique.

\end{document}
